% 第 13 题(3) 解析：四个正方形区域 Ⅰ–Ⅳ 与同轴控束（按原图矢量重绘）
\begin{tikzpicture}[font=\small, >=Stealth, scale=0.9]
  % 坐标轴
  \draw[->] (-5.2,0) -- (5.6,0) node[right] {$x$};
  \draw[->] (0,-3.6) -- (0,3.6) node[above] {$y$};
  \node[below left] at (0,0) {$O$};
  % 区域 Ⅰ（左上，边长 r3，磁场向里）
  \draw[thick,dashed] (-3.0,0.1) rectangle (-0.1,3.0);
  \node at (-1.55,2.5) {Ⅰ};
  % 区域 Ⅱ（左下，边长 r3，磁场向外）
  \draw[thick,dashed] (-3.0,-0.1) rectangle (-0.1,-3.0);
  \node at (-1.55,-2.5) {Ⅱ};
  % 区域 Ⅲ（右上，边长 r4）
  \draw[thick,dashed] (0.1,0.1) rectangle (3.6,3.6);
  \node at (3.1,3.1) {Ⅲ};
  % 区域 Ⅳ（右下，边长 r4）
  \draw[thick,dashed] (0.1,-0.1) rectangle (3.6,-3.6);
  \node at (3.1,-3.1) {Ⅳ};
  % 入射粒子流（宽 2r3）
  \draw[->,thick] (-5.0,3.0) -- (-3.05,3.0);
  \draw[->,thick] (-5.0,0.1) -- (-3.05,0.1);
  \draw[<->] (-4.7,0.1) -- (-4.7,3.0);
  \node[left,font=\scriptsize] at (-4.7,1.55) {$2r_{3}$};
  % 各区磁场（Ⅰ、Ⅲ 向里；Ⅱ、Ⅳ 向外）
  \foreach \x in {-2.5,-1.9,-1.3,-0.7} {
    \foreach \y in {0.6,1.2,1.8,2.4} {
      \draw (\x,\y) ++(-0.06,-0.06) -- ++(0.12,0.12);
      \draw (\x,\y) ++(-0.06,0.06) -- ++(0.12,-0.12);
    }
  }
  \foreach \x in {-2.5,-1.9,-1.3,-0.7} {
    \foreach \y in {-2.4,-1.8,-1.2,-0.6} {
      \fill (\x,\y) circle (1.3pt);
    }
  }
  \foreach \x in {0.6,1.4,2.2,3.0} {
    \foreach \y in {0.6,1.4,2.2,3.0} {
      \fill (\x,\y) circle (1.3pt);
    }
  }
  \foreach \x in {0.6,1.4,2.2,3.0} {
    \foreach \y in {-3.0,-2.2,-1.4,-0.6} {
      \draw (\x,\y) ++(-0.06,-0.06) -- ++(0.12,0.12);
      \draw (\x,\y) ++(-0.06,0.06) -- ++(0.12,-0.12);
    }
  }
  % 出射粒子流（宽 2r4）
  \draw[->,thick] (3.65,3.6) -- (5.4,3.6);
  \draw[->,thick] (3.65,0.1) -- (5.4,0.1);
  \draw[<->] (5.1,0.1) -- (5.1,3.6);
  \node[right,font=\scriptsize] at (5.15,1.85) {$2r_{4}$};
\end{tikzpicture}
